% =============================================================================
% Section: Baseline Models
% Baseline A (Classical ML) and Baseline B (Frozen Pretrained Encoder)
% SemEval-2027 Task 3 -- Swedish development set, Subtask 1 and Subtask 2.
% =============================================================================

\section{Baseline Models}

\subsection{Task Formulation}

SemEval-2027 Task~3 frames diachronic semantic change as two complementary
sub-problems over Swedish contextual usage records spanning 10~time periods
(five historical: 1880--1883, 1888--1891, 1896--1899, 1904--1907, 1912--1915;
five modern: 2006--2007, 2010--2011, 2014--2015, 2018--2019, 2022--2023).

\paragraph{Subtask~1 --- Diachronic Word Sense Induction (WSI).}
Given all contextual usages of a target word $w$ across all periods,
the system must assign each usage a sense label drawn from an induced sense
inventory.  No predefined sense hierarchy is provided; the task is
\emph{unsupervised}.  Gold labels are annotator-assigned sense identifiers.
A non-trivial fraction (13.2\,\% of development records, $n=262$ out of 1983)
carry \emph{multi-valued} annotations indicating genuine polysemy or
inter-annotator disagreement.  These records are preserved throughout but
excluded from ARI, NMI, and purity evaluation per the official multi-label
policy.

\paragraph{Subtask~2 --- Binary Usage Classification.}
For each contextual usage, the system predicts a binary label $y \in \{0,1\}$,
where label~1 indicates that the usage instantiates the target sense as defined
in the task's operational sense definitions.  This is a \emph{supervised}
binary classification task.  The development set contains 1983~records with
class distribution 0:1278 / 1:705, yielding an imbalance ratio of
approximately 1.81:1.

These two sub-problems are methodologically distinct.
Subtask~1 requires bottom-up sense discovery (clustering), whereas Subtask~2
requires top-down sense discrimination (classification) relative to a fixed
reference sense.  A system that performs well on Subtask~2 need not produce
interpretable sense clusters, and vice versa.


% -----------------------------------------------------------------------------
\subsection{Data Split and Leakage Prevention}
\label{sec:baselines:split}

The development set is split into a stratified 80/20 train--test partition
($\text{train}=1586$, $\text{test}=397$) using \texttt{train\_test\_split}
with \texttt{stratify=label\_value} and \texttt{random\_state}~$=42$.

Leakage prevention is enforced at the architectural level:
\begin{itemize}
  \item All vectorisers (\texttt{TfidfVectorizer}), scalers
    (\texttt{StandardScaler}), and categorical encoders (\texttt{OneHotEncoder})
    are encapsulated within \texttt{sklearn.pipeline.Pipeline} objects and
    fitted exclusively on the training split.
  \item For Baseline~B, frozen XLM-R embeddings are generated without label
    access.  Downstream classifiers receive only training-split embeddings
    during fitting.
  \item Deterministic seed~42 is fixed across \texttt{random}, \texttt{numpy},
    and \texttt{sklearn} for full reproducibility.
\end{itemize}

Subtask~1 clustering is \emph{unsupervised} and does not interact with the
train--test split.  Evaluation metrics are computed post-hoc against gold
singleton labels on the entire word's usage set.


% -----------------------------------------------------------------------------
\subsection{Baseline~A: Classical ML with Interpretable Features}

\subsubsection{Feature Design and Connection to EDA}
\label{sec:baselines:features}

The EDA phase established three findings that directly inform feature selection:
\begin{enumerate}
  \item \textbf{Per-word label imbalance.} The label-1 rate is highly
    word-dependent: \emph{exportera} (94.6\,\%) and \emph{fru} (89.4\,\%)
    contrast sharply with \emph{fr\"{o}ken} (5.7\,\%) and
    \emph{herre} (7.1\,\%).  A model with access to target-word identity
    can exploit these base rates without genuine semantic reasoning.
  \item \textbf{Structural position signal.} The character-offset position of
    the target token (\texttt{target\_position}) was identified as a
    discriminative structural covariate in the EDA analysis.
  \item \textbf{Variable lexical separability.} EDA embedding projections
    showed partial separability for \emph{fru} and \emph{kvinna} but
    near-random structure for \emph{fr\"{o}ken} and \emph{exportera}.
\end{enumerate}

The full interpretable feature matrix combines three channels:

\paragraph{Lexical channel (TF-IDF).}
\texttt{TfidfVectorizer} over word unigrams and bigrams,
\texttt{sublinear\_tf=True}, \texttt{min\_df=1}, vocabulary cap 50\,000.
For the Linear~SVM, a character-level TF-IDF (3--5-grams) is concatenated
via \texttt{FeatureUnion} to capture morphological patterns.

\paragraph{Numeric EDA channel.}
Nine features from the EDA pipeline:
\texttt{text\_length\_chars}, \texttt{token\_count},
\texttt{unique\_token\_count}, \texttt{type\_token\_ratio},
\texttt{avg\_token\_length}, \texttt{punctuation\_ratio},
\texttt{digit\_ratio}, \texttt{stopword\_ratio},
\texttt{target\_position}. All standardised with \texttt{StandardScaler}.

\paragraph{Categorical channel.}
One-hot encodings of \texttt{word} (target-word identity),
\texttt{pos} (part-of-speech), and \texttt{period\_label},
enabling per-word and per-period intercepts.

\subsubsection{Word-Identity Ablation}
\label{sec:baselines:ablation}

A dedicated ablation condition isolates contextual signal from token-identity
memorisation:
\begin{itemize}
  \item The target token is replaced with \texttt{[TARGET]} via case-insensitive,
    word-boundary-delimited regex substitution before vectorisation.
  \item The categorical \texttt{word} feature is excluded from the transformer.
\end{itemize}
This model measures how much classification signal is attributable to
surrounding context and structure alone.

\subsubsection{TF-IDF Model Configurations}

Seven classical models are evaluated for Subtask~2:
\begin{enumerate}
  \item \textbf{Majority baseline} (\texttt{DummyClassifier},
    \texttt{most\_frequent}): performance floor.
  \item \textbf{TF-IDF + LR}: word 1,2-grams; \texttt{class\_weight='balanced'}.
  \item \textbf{Word+Char TF-IDF + LinearSVC}: word 1,2-grams and
    char 3--5-grams; \texttt{class\_weight='balanced'}.
  \item \textbf{TF-IDF + Complement NB} (\texttt{norm=True}): complement-class
    weighting for class imbalance.
  \item \textbf{EDA-feature LR}: TF-IDF + numeric EDA + categorical; LR.
  \item \textbf{EDA-feature RF}: 100 estimators, \texttt{max\_depth=15},
    \texttt{class\_weight='balanced'}.
  \item \textbf{Ablated LR}: word masked, \texttt{word} categorical excluded.
\end{enumerate}

\subsubsection{Interpretability: LR Coefficients and RF Gini Importances}
\label{sec:baselines:interpret}

Linear model coefficients quantify the log-odds contribution of each feature.
For the TF-IDF LR, top positive coefficients (label~1 association) are
\emph{fru} ($+4.37$), \emph{exportera} ($+2.02$), \emph{sin fru} ($+1.60$),
\emph{exporterar} ($+1.57$).  Top negative coefficients (label~0) are
\emph{fr\"{o}ken} ($-2.12$), \emph{herrar} ($-1.46$), \emph{motivet}
($-1.17$).

These patterns reflect per-word label-prevalence, not semantic change
detection.  The token \emph{fru} carries a large positive coefficient
because 89.4\,\% of \emph{fru} usages are labelled~1, confirming EDA
finding~(1): the model exploits target-word vocabulary as a class-rate proxy.

For the EDA-feature LR, \texttt{categorical\_\_word\_fru} ($+2.98$) and
\texttt{categorical\_\_word\_exportera} ($+2.79$) dominate, with
\texttt{numeric\_\_text\_length\_chars} ($+0.62$) in the top~15,
corroborating EDA finding~(2).

\paragraph{Random Forest Gini Importances.}
Among 45,190~features, the top three by mean Gini decrease are:
\texttt{categorical\_\_word\_exportera} (0.0248),
\texttt{categorical\_\_word\_herre} (0.0204), and
\texttt{numeric\_\_target\_position} (0.0199, ranked \#3 overall).
The high ranking of \texttt{target\_position} validates EDA finding~(2).

\subsubsection{Subtask~1: Diachronic Sense Induction via Clustering}

TF-IDF KMeans ($k=3$) is applied \emph{per target word}, augmented with
nine standardised numeric EDA features concatenated to the sparse TF-IDF
matrix.  Clustering uses the entire word's usage set (no train--test split).

Critical interpretive constraints:
\begin{itemize}
  \item Cluster assignments are \textbf{exploratory groupings} by co-text
    similarity, not validated sense labels.
  \item A dominant-cluster shift across periods may equally reflect
    \emph{corpus source}, \emph{genre}, or \emph{topical vocabulary} change.
    No causal claim of semantic change is made without further validation.
  \item High cluster purity can arise from skewed gold-label distributions
    rather than genuine sense boundary discovery.
\end{itemize}


% -----------------------------------------------------------------------------
\subsection{Baseline~B: Frozen Pretrained Encoder (XLM-R)}

\subsubsection{Feature Extraction vs.\ Fine-Tuning}

Baseline~B uses \textbf{frozen representations} from
\texttt{FacebookAI/xlm-roberta-base}.  Encoder parameters are not updated.
This is methodologically distinct from fine-tuning:
\begin{itemize}
  \item \textbf{Feature extraction} (this approach): encoder is a fixed
    embedding function.  Downstream classifiers receive mean-pooled hidden
    states.  Representation quality is determined solely by the pretrained
    model's coverage of Swedish diachronic text.
  \item \textbf{Fine-tuning}: encoder parameters are updated on task data,
    improving task-specific fit at the cost of interpretability and with
    overfitting risk on the small development set.
\end{itemize}
Frozen-encoder results establish a reproducible reference point against which
fine-tuned models can be compared.

\subsubsection{Contextual Embedding Method}

Input to the encoder is a $\pm80$-character window around the target token.
Mean-pooling over the final hidden-layer token representations (masked by the
attention mask) produces a 768-dimensional vector per usage.  Batch size: 32;
maximum token length: 512.

Embeddings are SHA-256 fingerprinted over all usage records and context strings
and cached on disk, with validation of model name, pooling method, sequence
length, and data fingerprint before reuse.

\subsubsection{Downstream Classification and Clustering}

\paragraph{Subtask~2.}
Each 768-d embedding is standardised (\texttt{StandardScaler}) and passed to
Logistic Regression and Linear~SVM (both \texttt{class\_weight='balanced'}),
isolating representational quality without non-linear capacity confounds.

\paragraph{Subtask~1.}
Per-word KMeans ($k=3$, \texttt{n\_init=10}) is applied directly on the
embedding matrix without additional feature augmentation.


% -----------------------------------------------------------------------------
\subsection{Evaluation Metrics}

\paragraph{Subtask~2.}
Primary metric: \textbf{macro-averaged F1} (F1-macro), giving equal weight
to each binary class.  Raw accuracy is unsuitable as primary metric: the
majority-class predictor achieves 64.5\,\% while failing entirely on the
minority class.  Balanced accuracy (arithmetic mean of per-class recall)
is the secondary imbalance-aware metric.

\paragraph{Subtask~1.}
\begin{itemize}
  \item \textbf{ARI}: cluster-vs-gold agreement corrected for chance.
    ARI~$=0$ is chance-level; ARI~$<0$ is worse-than-chance.
  \item \textbf{NMI}: shared information between cluster and gold-label
    distributions, normalised to $[0,1]$.
  \item \textbf{Purity}: proportion of each cluster dominated by one gold
    sense.  High purity is necessary but not sufficient for good induction.
\end{itemize}
All Subtask~1 metrics are computed exclusively on singleton-labelled records;
multi-label records are preserved but excluded from metric computation.


% -----------------------------------------------------------------------------
\subsection{Results}

\subsubsection{Subtask~2: Binary Usage Classification}

\begin{table}[H]
\centering
\caption{Subtask~2 results on the Swedish development test split ($n=397$,
  stratified 20\,\%).  F1-macro is the primary ranking metric.
  BA\,=\,balanced accuracy.  Sorted by F1-macro descending.}
\label{tab:subtask2}
\small
\begin{tabular}{llrrr}
\toprule
\textbf{Model} & \textbf{Representation} & \textbf{F1-macro} & \textbf{BA} & \textbf{Acc.} \\
\midrule
EDA-feat.\ LR        & TF-IDF + EDA features    & \textbf{0.8163} & \textbf{0.8049} & 83.9\,\% \\
Word+Char TF-IDF SVM & Word \& char TF-IDF      & 0.8155          & 0.8033          & 83.9\,\% \\
TF-IDF LR            & Word TF-IDF              & 0.7917          & 0.7826          & 81.6\,\% \\
EDA-feat.\ RF        & TF-IDF + EDA features    & 0.7910          & 0.7741          & 82.4\,\% \\
Ablated LR           & Context only, no word ID & 0.7129          & 0.7165          & 73.3\,\% \\
Complement NB        & Word TF-IDF              & 0.6171          & 0.6253          & 73.0\,\% \\
Majority baseline    & Constant label\,=\,0     & 0.3920          & 0.5000          & 64.5\,\% \\
\midrule
Frozen XLM-R LR      & 768-d mean-pooled XLM-R  & \multicolumn{3}{c}{\emph{pending (model weights downloading)}} \\
Frozen XLM-R SVM     & 768-d mean-pooled XLM-R  & \multicolumn{3}{c}{\emph{pending (model weights downloading)}} \\
\bottomrule
\end{tabular}
\end{table}

The best classical model, EDA-feature LR (F1-macro\,$=0.8163$), achieves
$+42.4$~points over the majority baseline.  The Word+Char TF-IDF SVM is
essentially tied ($0.8155$), indicating that character-level morphological
features provide only marginal additional signal over word-level features.

Complement~NB ($0.6171$) achieves $22.5$~points below the TF-IDF LR ($0.7917$)
despite the same representation.  This reflects an unfavourable interaction
between the complement-class weighting assumption and this vocabulary size
and class distribution.

\subsubsection{Subtask~1: Diachronic Word Sense Induction}

\begin{table}[H]
\centering
\caption{Subtask~1 clustering metrics per target word (TF-IDF KMeans,
  $k=3$, singleton labels only). Records: total usages; Eval.~$n$: singleton
  records used for metrics; Excl.: multi-label records excluded.}
\label{tab:subtask1}
\small
\begin{tabular}{lrrrrrr}
\toprule
\textbf{Word} & \textbf{Rec.} & \textbf{Eval.~$n$} & \textbf{Excl.} & \textbf{ARI} & \textbf{NMI} & \textbf{Purity} \\
\midrule
\emph{fru}        & 274 & 195 & 79 & \textbf{0.255} & 0.096 & 0.867 \\
\emph{motiv}      & 291 & 289 &  2 & 0.144          & 0.051 & 0.785 \\
\emph{kvinna}     & 290 & 284 &  6 & 0.060          & 0.039 & 0.775 \\
\emph{dumpa}      & 145 & 103 & 42 & 0.042          & 0.107 & 0.815 \\
\emph{herre}      & 283 & 241 & 42 & 0.040          & 0.066 & 0.838 \\
\emph{suga}       & 129 & 127 &  2 & 0.038          & 0.082 & 0.323 \\
\emph{exportera}  & 149 & 149 &  0 & 0.005          & 0.042 & 0.946 \\
\emph{sk\"{a}r}  &  98 &  95 &  3 & $-$0.003       & 0.052 & 0.547 \\
\emph{f\"{o}rort} & 150 &  71 & 79 & $-$0.015       & 0.018 & 0.577 \\
\emph{fr\"{o}ken} & 174 & 167 &  7 & $-$0.104       & 0.065 & 0.838 \\
\bottomrule
\end{tabular}
\end{table}

ARI values are near zero for the majority of words and negative for three
(\emph{fr\"{o}ken}: $-0.104$; \emph{f\"{o}rort}: $-0.015$;
\emph{sk\"{a}r}: $-0.003$).  Highest ARI is achieved by \emph{fru} ($0.255$).
\emph{Exportera} presents the highest purity ($0.946$) despite near-zero ARI,
consistent with a single dominant sense accounting for most usages.

Temporal dominant-cluster tracking shows 6 of 10~words exhibit at least
two cluster transitions (\emph{dumpa}, \emph{exportera}, \emph{fru},
\emph{f\"{o}rort}, \emph{herre}, \emph{kvinna}).  Four words
(\emph{fr\"{o}ken}, \emph{motiv}, \emph{sk\"{a}r}, \emph{suga}) show no
dominant-cluster shift.  These patterns are consistent with temporal lexical
variation but \textbf{must not} be interpreted as evidence of semantic change
without further validation.


% -----------------------------------------------------------------------------
\subsection{Error Analysis by Target Word and Period}
\label{sec:baselines:error}

\subsubsection{Per-Word Error Analysis}

Per-word group metrics for the best classical model (EDA-feature LR) show
highly uneven per-word performance:
\begin{itemize}
  \item \emph{Fru}: 94.3\,\% accuracy but F1-macro\,$\approx0.69$.
    High accuracy reflects the 89.4\,\% label-1 base rate.
  \item \emph{Fr\"{o}ken} and \emph{herre}: $\approx$92--93\,\% accuracy
    but F1-macro\,$\approx0.48$, near majority-baseline performance, due to
    94.3\,\% and 92.9\,\% label-0 base rates respectively.
  \item \emph{Kvinna}: weakest per-word F1-macro ($\approx0.39$--$0.47$),
    confirming its usage patterns are the most lexically ambiguous.
  \item \emph{Dumpa} and \emph{suga}: intermediate F1-macro ($0.52$--$0.65$).
\end{itemize}

The ablated model (F1-macro\,$=0.7129$) maintains this structure without
target-word identity.  The $\Delta\text{F1}\approx10.3$~points gap to the
best full model ($0.8163$) quantifies the token-identity shortcut contribution.

\subsubsection{Per-Period Error Analysis}

Modern periods (2006--2023) generally yield higher F1-macro than historical
periods (1880--1915).  The TF-IDF~LR achieves F1-macro\,$=0.927$ on the
2022--2023 test subset versus $0.689$ on 1880--1883.  This gap has several
confounding causes:
\begin{itemize}
  \item Modern vocabulary is better represented in TF-IDF training vocabulary.
  \item Historical text contains archaic orthography and syntactic patterns
    that are under-represented in training data.
  \item Within-period class ratios vary and are not controlled separately
    from the overall stratified split.
\end{itemize}
This gap therefore cannot be interpreted as the model detecting semantic change.


% -----------------------------------------------------------------------------
\subsection{Exploratory vs.\ Predictive Evidence and Causal Caution}
\label{sec:baselines:caution}

A central methodological requirement in diachronic NLP is the distinction
between exploratory and predictive evidence:

\begin{description}
  \item[Exploratory evidence] (EDA and Subtask~1 clustering) identifies
    patterns that may correspond to meaningful linguistic phenomena but have
    not been subjected to held-out validation.  Cluster separability,
    embedding neighbourhood structure, and temporal distribution shifts
    motivate hypotheses and guide feature design, but do not establish
    causal or predictive claims.
  \item[Predictive evidence] (Subtask~2 classification) is established by
    measuring held-out generalisation on the test split.  Improvements in
    F1-macro relative to the majority baseline demonstrate that certain features
    carry signal about the binary label.  However, the binary label is defined
    operationally by annotators, and high F1 may reflect annotator conventions
    rather than fundamental semantic change.
\end{description}

\paragraph{Semantic change vs.\ corpus composition effects.}
Neither embedding-space cluster separation nor temporal cluster shifts
constitute \emph{direct} evidence of semantic change.  Alternative explanations
that must be systematically ruled out include:
\begin{itemize}
  \item \textbf{Genre and register shift.} Historical newspaper text differs
    from modern web text in syntactic complexity, topic distribution, and
    discourse structure.  Cluster shifts may reflect genre rather than
    lexical meaning change.
  \item \textbf{Source corpus vocabulary change.} Different source corpora
    may dominate different periods, introducing vocabulary shifts that TF-IDF
    captures as cluster transitions even if the word's sense is stable.
  \item \textbf{Sample imbalance.} Ten target words with 98--291~records each
    is a small and potentially non-representative sample.
  \item \textbf{Annotator conventions.} Words with extreme label rates
    (\emph{exportera}: 94.6\,\% label~1) may primarily reflect how annotators
    applied the operational sense definition.
\end{itemize}


% -----------------------------------------------------------------------------
\subsection{Limitations}

\begin{enumerate}
  \item \textbf{Single development set, no cross-validation.}
    Standard deviation estimates are unavailable; results may be sensitive
    to the specific 80/20 partition.
  \item \textbf{Fixed cluster count.}
    $k=3$ is applied uniformly; automatic selection (silhouette, elbow) is
    a natural next step.
  \item \textbf{Bag-of-words limitations.}
    TF-IDF discards word order and syntactic structure, which are potentially
    important for sense discrimination in historical text.
  \item \textbf{Frozen encoder domain mismatch.}
    XLM-R was pretrained on contemporary web text.  Its representations of
    19th-century Swedish may be limited.  Embedding distance does not prove
    semantic change; it reflects the encoder's modern corpus prior.
  \item \textbf{Multi-label exclusions.}
    The 13.2\,\% multi-label records excluded from Subtask~1 metrics are
    likely the most linguistically interesting borderline usages, and
    excluding them may yield optimistic estimates.
  \item \textbf{Class imbalance handling.}
    \texttt{class\_weight='balanced'} is the only mitigation strategy used.
    Threshold calibration and oversampling are not explored.
  \item \textbf{Baseline~B pending.}
    Frozen XLM-R results are not yet available due to model-weight download
    constraints.  The evaluation code is in place and will generate results
    automatically once the cache is complete.
\end{enumerate}
